\documentclass[11pt,a4paper]{article}
\usepackage[utf8]{inputenc}
\usepackage[T1]{fontenc}
\usepackage{CJKutf8}
\usepackage{geometry}
\geometry{margin=2.5cm}
\usepackage{booktabs}
\usepackage{amsmath,amssymb}
\usepackage{graphicx}
\usepackage{xcolor}
\usepackage{enumitem}
\usepackage{hyperref}
\usepackage{fancyhdr}
\usepackage{titlesec}

\pagestyle{fancy}
\fancyhf{}
\rhead{2026-05-13}
\lhead{SRS 2D MMSE Progress Report}
\rfoot{\thepage}

\titleformat{\section}{\large\bfseries}{}{0em}{}
\titleformat{\subsection}{\normalsize\bfseries}{}{0em}{}

\begin{document}
\begin{CJK}{UTF8}{gbsn}

\begin{center}
{\LARGE\bfseries SRS 2D MMSE 信道估计 进展报告}\\[0.8em]
{}
\end{center}

\vspace{1em}

%% ================================================================
\section{一、教授指导与工作思路}

\subsection{1.1 教授 4/29 指导}

在 OAI SRS 信道估计中引入 2D MMSE filtering，替代当前的 LS + 固定 FIR（filt8/16）pipeline。

\subsection{1.2 教授 5/7 组会反馈}

\begin{enumerate}[leftmargin=2em]
\item \textbf{PDP-MMSE}：可以先从 LS 估计中获取 PDP，再确定 MMSE 系数；也可以根据 PDP / Doppler 等级预设一组滤波器
\item \textbf{自适应滤波 trade-off}：LMS 收敛快但性能差，RLS 性能好但复杂，Lattice-RLS 取平衡
\item \textbf{复杂度是硬约束}：SRS 块级别的执行时间增加必须 negligible
\end{enumerate}

\subsection{1.3 本周工作思路}

\begin{enumerate}[leftmargin=2em]
\item 探索三种候选方案（频域 MMSE、时域 EWMA、Lattice-RLS），评估性能-复杂度 trade-off
\item 发现并诊断 Legacy pipeline 的 NMSE floor（$\approx -7$\,dB），定位根因至 int16 信号幅度
\item 确定 2D MMSE 的正确架构：可分离滤波（1D 时域 Wiener + 1D 频域 Wiener）
\item 修复了 Proxy 信道模型中一个导致"伪静态信道"的 bug
\end{enumerate}

%% ================================================================
\section{二、方案探索与选择}

\subsection{2.1 方案 A：频域 PDP$\to$R MMSE}

实现了完整的 1D 频域 MMSE 信道估计器：
\[
\text{LS} \;\to\; \text{IFFT} \;\to\; \text{PDP} \;\to\; R(\Delta k) \;\to\; \text{per-SC MMSE solve}
\]
在 OAI \texttt{nr\_srs\_mmse.c} 中新增约 160 行代码。

\begin{table}[h]\centering
\begin{tabular}{lcc}
\toprule
指标 & Legacy & PDP$\to$R MMSE \\
\midrule
NMSE p50（SNR=20\,dB） & $-7.12$\,dB & $-2.46$\,dB \\
单帧执行时间 & $\sim$0.1\,ms & \textbf{5--10\,ms} \\
\bottomrule
\end{tabular}
\end{table}

\textbf{淘汰原因}：处理时间增加 50--100$\times$，不满足 negligible 约束。但验证了频域 MMSE 替代 filt8/16 的\textbf{方向正确}。

\subsection{2.2 方案 B：时域 EWMA 平滑——已验证有效}

实现了指数加权移动平均（EWMA）时域滤波：
\[
\hat{H}_{\text{smooth}}(t) = \alpha \cdot \hat{H}_{\text{smooth}}(t\!-\!1) + (1\!-\!\alpha) \cdot \hat{H}_{\text{new}}(t)
\]

\textbf{静态信道实测}（speed=0, SNR=20\,dB, seed=42）：

\begin{table}[h]\centering
\begin{tabular}{lccc}
\toprule
指标 & Legacy & EWMA $\alpha=0.9$ & 改善 \\
\midrule
NMSE p50 & $-6.53$\,dB & $\mathbf{-7.08}$\,dB & $+0.55$\,dB \\
NMSE p90 & $-0.89$\,dB & $\mathbf{-4.41}$\,dB & $+3.52$\,dB \\
spread (p90$-$p10) & 6.68\,dB & 3.55\,dB & 缩小 3.13\,dB \\
额外执行时间 & --- & negligible & --- \\
\bottomrule
\end{tabular}
\end{table}

EWMA 验证了时域平滑对 SRS 信道估计的有效性（2D MMSE 的时间维度）。

\subsection{2.3 方案 C：Lattice-RLS M=2}

参照 Haykin Ch.16 实现了 M=2 的 LSL 并做了验证。

\textbf{结论}：Lattice-RLS 在 adaptive equalization 场景下（$d = \text{known pilot},\; u = \text{received}$）有效。但在 channel state smoothing 场景下（$d = u = \text{noisy LS}$），joint-process estimator 的最优解退化为 identity filter（$\rho_0=1$, 其余为 0），这是数学性质而非实现 bug。

\subsection{2.4 方案选择小结}

\begin{table}[h]\centering
\begin{tabular}{lccc}
\toprule
方案 & NMSE 改善 & 复杂度增加 & 结论 \\
\midrule
PDP$\to$R 频域 MMSE & 方向对，实现过重 & 50--100$\times$ & 淘汰 \\
\textbf{EWMA 时域平滑} & \textbf{p50 +0.55, p90 +3.52\,dB} & \textbf{negligible} & \textbf{已验证} \\
Lattice-RLS & 退化为 identity & 中等 & 不适用 \\
\bottomrule
\end{tabular}
\end{table}

%% ================================================================
\section{三、$-4.8$\,dB 单帧精度瓶颈：诊断与分析}

在方案探索过程中，我们发现无论使用何种算法（Legacy LS+filt8 / PDP$\to$R MMSE / EWMA），OAI 的单帧 SRS 信道估计 NMSE 始终无法突破约 $-7$\,dB。这个 floor 既不随 SNR 改善（SNR 从 20 提高到 40\,dB 仅改善 0.36\,dB），也不随插值方法改变（跳过 filt8 仅改善 0.15\,dB）。为此进行了系统性诊断。

\subsection{3.1 控制变量实验}

为定位 floor 来源，我们设计了一组控制变量实验，逐一排除可能的因素：

\begin{table}[h]\centering
\begin{tabular}{p{4.5cm}p{3.5cm}p{5.5cm}}
\toprule
实验 & 目的 & 结果 \\
\midrule
passthru（跳过 filt8，仅输出 LS） & 验证 filt8 是否是 floor 的主因 & NMSE 仅差 0.15\,dB $\to$ \textbf{filt8 不是主因} \\
\addlinespace
静态信道（speed=0） & 消除 Doppler 和帧间对齐干扰 & floor 依然 $-6.5$\,dB $\to$ \textbf{与移动性无关} \\
\addlinespace
SNR=40\,dB 静态信道 & 消除 thermal noise 的贡献 & p50 仅改善 0.36\,dB $\to$ \textbf{与 SNR 无关} \\
\addlinespace
Float64 仿真（Python，filt8/sinc/Wiener） & 排除插值算法和系数设计问题 & float 下所有方法达 $-20\sim -22$\,dB $\to$ \textbf{int16 链路才是差距来源} \\
\bottomrule
\end{tabular}
\end{table}

\textbf{关键推理}：
\begin{itemize}[leftmargin=2em]
\item passthru $\approx$ Legacy（差 0.15\,dB）：filt8 既不帮忙也不添乱，floor 在 filt8 \textbf{上游}（LS 估计本身）
\item Float64 仿真达 $-22$\,dB：同样的 filt8 算法在 float 精度下表现正常，说明 floor 来自 \textbf{int16 定点精度}
\item SNR=40 几乎无改善：floor 不是 thermal noise，而是 int16 链路引入的\textbf{量化失真}
\end{itemize}

\subsection{3.2 根因定位：int16 信号幅度过低}

进一步追踪 int16 链路各级的信号幅度，发现了根本原因：

OAI 的 SRS 参考信号幅度 AMP=512（\texttt{impl\_defs\_top.h} 中定义），经过 OFDM 调制（IFFT）后，时域信号的 RMS 仅约 200，占 int16 范围（$\pm$32767）的 \textbf{0.6\%}——即 \textbf{只有 7.6 bit 的有效精度}（int16 共 15 bit）。

\begin{table}[h]\centering
\begin{tabular}{lccc}
\toprule
处理级别 & 信号 RMS & 占 int16 范围 & 有效位数 \\
\midrule
UE IFFT 时域输出 & 200 & \textbf{0.6\%} & 7.6\,bit \\
Proxy$\to$shm 量化后 & $\sim$564 & 1.7\% & 9.1\,bit \\
OAI FFT 频域输出 & 14000 & 42.7\% & 13.8\,bit \\
LS 除法输出 & 9900 & 30.2\% & 13.3\,bit \\
\bottomrule
\end{tabular}
\end{table}

\textbf{关键实测数据}（静态信道 speed=0, SNR=20\,dB, 195 帧）：

\begin{table}[h]\centering
\begin{tabular}{lcc}
\toprule
指标 & 预期值 & 实测值 \\
\midrule
SRS 帧间 correlation & 0.99 & \textbf{0.80} \\
Per-SC 有效 SNR (mean) & 20\,dB & \textbf{7.0\,dB} \\
Per-SC 有效 SNR (p10/p50/p90) & --- & $-1.7$ / $5.1$ / $19.5$\,dB \\
单帧 NMSE (vs 50帧平均 GT) & $-20$\,dB & $\mathbf{-4.8}$\,dB \\
\bottomrule
\end{tabular}
\end{table}

\textbf{小结}：在配置 SNR=20\,dB 的条件下，单帧 LS 的 NMSE 仅 $-4.8$\,dB（有效 SNR 约 7\,dB，损失 13\,dB）。这说明 OAI 的单帧 LS 估计精度受限于 int16 信号幅度，而非算法或 thermal noise。

\subsection{3.3 验证：时域多帧平均可以突破 floor}

既然 floor 来自 per-frame 的 int16 量化失真（帧间近似独立），多帧平均应当以 $+10\log_{10}(N)$\,dB 的速率降低 NMSE。我们在静态信道数据上验证了这一点：

\begin{table}[h]\centering
\begin{tabular}{rrrrrr}
\toprule
$W$ & NMSE (all) & NMSE (pilot) & NMSE (mid) & 实测 gain & 理论 $10\lg W$ \\
\midrule
1 & $-4.8$\,dB & $-4.8$\,dB & $-4.8$\,dB & 0\,dB & 0\,dB \\
2 & $-9.1$ & $-9.1$ & $-9.0$ & $+4.3$ & $+3.0$ \\
5 & $-12.8$ & $-12.8$ & $-12.7$ & $+7.9$ & $+7.0$ \\
10 & $-14.8$ & $-14.8$ & $-14.7$ & $+10.0$ & $+10.0$ \\
20 & $-16.9$ & $-16.9$ & $-16.9$ & $+12.1$ & $+13.0$ \\
50 & $-20.8$ & $-20.9$ & $-20.8$ & $+16.0$ & $+17.0$ \\
100 & $-24.3$ & $-24.3$ & $-24.3$ & $+19.5$ & $+20.0$ \\
\bottomrule
\end{tabular}
\end{table}

\textbf{观察}：
\begin{enumerate}[leftmargin=2em]
\item 实测增益几乎完美追踪理论值 $10\log_{10}(W)$，在 $W=10$ 时达到 $+10.0$\,dB（理论 $+10.0$）
\item Pilot SC 与 Midpoint SC 的 NMSE 完全一致，再次确认 filt8 不引入额外误差
\item $W=100$ 帧时 NMSE 达 $-24.3$\,dB，说明 floor 并非不可突破
\end{enumerate}

\textbf{EWMA 验证}（与已实现的 OAI EWMA 对标）：

\begin{table}[h]\centering
\begin{tabular}{rrrr}
\toprule
$\alpha$ & $N_{\text{eff}} = 1/(1\!-\!\alpha)$ & NMSE & Gain \\
\midrule
0.00 & 1 & $-4.8$\,dB & 0\,dB \\
0.50 & 2 & $-10.5$ & $+5.7$ \\
0.80 & 5 & $-14.6$ & $+9.8$ \\
0.90 & 10 & $-16.8$ & $+11.9$ \\
0.95 & 20 & $-18.6$ & $+13.8$ \\
0.98 & 50 & $-20.2$ & $+15.4$ \\
\bottomrule
\end{tabular}
\end{table}

\subsection{3.4 对 2D MMSE 设计的指导意义}

Floor 诊断直接决定了 2D MMSE 的正确架构：

\begin{table}[h]\centering
\begin{tabular}{p{5cm}p{8cm}}
\toprule
诊断结论 & 对 2D MMSE 设计的影响 \\
\midrule
passthru $\approx$ legacy（filt8 不是瓶颈） & 不必保留 filt8，可以用频域 MMSE 完全替代 \\
\addlinespace
单帧精度 $-4.8$\,dB 受限于 int16 幅度 & \textbf{时域多帧平均是提升精度的主力}（$W=10$ 即可 $+10$\,dB） \\
\addlinespace
增益追踪 $10\log_{10}(W)$ 理论值 & 帧间量化失真近似独立，MMSE 时域滤波可获得最优增益 \\
\bottomrule
\end{tabular}
\end{table}

%% ================================================================
\section{四、下一步——可分离 2D MMSE}

\subsection{4.1 架构}

基于根因分析，确定了正确的 2D MMSE 架构（工业界标准的可分离滤波）：

\begin{center}
\fbox{\parbox{0.85\textwidth}{
\textbf{OAI 现有}：$\text{LS}_{624} \;\to\; \text{filt8}_{624\to 1248} \;\to\; \text{输出}$ \quad [单帧, 有效 SNR $\approx$ 2\,dB] \\[0.5em]
\textbf{2D MMSE}：$\text{LS}_{624} \;\to\; \underbrace{\text{时域 Wiener}}_{\text{跨 N 帧}} \;\to\; \underbrace{\text{频域 Wiener}}_{\text{替代 filt8, } 624\to 1248} \;\to\; \text{输出}$
}}
\end{center}

\textbf{关键设计}：时域 Wiener 在 filt8 \textbf{之前}（624 pilot SC 上）做，频域 Wiener \textbf{替代} filt8。

\textbf{原因}：filt8 将白噪声有色化，导致后续频域 MMSE 的公式
\[
\mathbf{w} = \left(\mathbf{R}_{HH} + \sigma^2 \mathbf{I}\right)^{-1} \mathbf{r}_{Hd}
\]
失效（$\sigma^2 \mathbf{I}$ 要求噪声独立）。在 pilot SC 上先做时域滤波，噪声保持白色，MMSE 最优性成立。

\subsection{4.2 复杂度与性能预期}

\begin{table}[h]\centering
\begin{tabular}{lccl}
\toprule
步骤 & 运算量 & 耗时 & 说明 \\
\midrule
时域 Wiener & $W \times 624 \times 4$ & $\sim$0.05\,ms & 在 pilot SC 上跨帧联合 \\
频域 Wiener & $624 \times 4 \times 4$ & $\sim$0.01\,ms & 替代 filt8, MMSE 插值 \\
\midrule
\textbf{总计} & --- & $\mathbf{\sim 0.06}$\,\textbf{ms} & 满足 negligible 要求 \\
\bottomrule
\end{tabular}
\end{table}

\textbf{内存}：$W=20$ 帧 $\times$ 2(rx) $\times$ 2(tx) $\times$ 624(pilot) $\times$ 4(bytes) = \textbf{200\,KB} 环形缓冲。

\textbf{预期性能}（静态信道 SNR=20\,dB）：

\begin{itemize}[leftmargin=2em]
\item 时域 $W=10$ 帧平均：有效 SNR 从 2.3 $\to$ \textbf{12.3\,dB}（$+10$\,dB）
\item 加频域 MMSE：额外 \textbf{3$\sim$5\,dB} 频域增益
\item 总计：NMSE 从 $-2$\,dB 改善到 $\mathbf{-15 \sim -17}$\,\textbf{dB}
\end{itemize}

\subsection{4.3 待完成}

\begin{itemize}[leftmargin=2em]
\item Python 离线原型验证完整 2D MMSE（时域 Wiener + 频域 Wiener）
\item OAI C 代码实现 + 多 SNR 点 sweep 实测
\item NMSE vs SNR 曲线 + 执行时间 profiling
\end{itemize}

%% ================================================================
\section{五、附：产出清单}

\subsection{OAI 代码}

\begin{table}[h]\centering
\begin{tabular}{ll}
\toprule
文件 & 内容 \\
\midrule
\texttt{nr\_srs\_mmse.c/h} & PDP$\to$R 频域 MMSE + estimator dispatcher \\
\texttt{nr\_srs\_2d\_filter.c/h} & EWMA 时域平滑（将扩展为完整 2D MMSE） \\
\texttt{nr\_ul\_channel\_estimation.c} & dispatcher 扩展 + SRS X\_ref dump \\
\bottomrule
\end{tabular}
\end{table}

\subsection{评估与诊断工具}

\begin{table}[h]\centering
\begin{tabular}{ll}
\toprule
文件 & 用途 \\
\midrule
\texttt{eval\_pdpR\_nmse.py} & NMSE 评估（LS-aligned + per-frame percentile） \\
\texttt{oai\_dft\_wrapper.py} & OAI bit-exact int16 DFT2048 Python 封装 \\
\texttt{test\_sinc\_interpolation.py} & Float64 插值方法对比仿真 \\
\texttt{preflight.sh} & sweep 前清理 + launch\_all 集成 \\
\texttt{run\_q4\_snr\_sweep\_v8.sh} & SNR sweep driver \\
\bottomrule
\end{tabular}
\end{table}

\end{CJK}
\end{document}
